\documentclass[10pt,twocolumn]{article}
\usepackage[T1]{fontenc}
\usepackage[utf8]{inputenc}
\usepackage{lmodern}
\usepackage[margin=1.8cm,top=2.4cm,bottom=2cm,columnsep=0.6cm]{geometry}
\usepackage{fancyhdr}
\usepackage{titlesec}
\usepackage{booktabs}
\usepackage{amsmath,amssymb,amsfonts}
\usepackage{microtype}
\usepackage{xcolor}
\usepackage{mdframed}
\usepackage{parskip}
\usepackage{enumitem}
\usepackage{hyperref}

\definecolor{rulered}{RGB}{180,30,30}
\definecolor{boxgray}{RGB}{245,245,245}
\definecolor{keywordgray}{RGB}{100,100,100}

\pagestyle{fancy}
\fancyhf{}
\fancyhead[L]{\textsc{\small Claude Mag}}
\fancyhead[C]{\small\textit{Machine Learning Research Digest}}
\fancyhead[R]{\small 2026-09-30}
\fancyfoot[C]{\small\thepage}
\renewcommand{\headrulewidth}{0.8pt}

\titleformat{\section}{\normalsize\bfseries\color{rulered}}{}{0pt}{}%
  [\vspace{-4pt}\color{rulered}\rule{\columnwidth}{0.4pt}]
\titlespacing{\section}{0pt}{8pt}{4pt}

\hypersetup{colorlinks=true,urlcolor=rulered,linkcolor=black}

\begin{document}
\setlength{\parindent}{0pt}
\setlength{\parskip}{4pt}

{\small\color{keywordgray}\textbf{Keywords:}
  floating-point precision \textbullet{}
  greedy decoding \textbullet{}
  LLM reproducibility \textbullet{}
  numerical divergence}\par\vspace{4pt}

{\LARGE\bfseries\raggedright
  Greedy Decoding Is Not Precision-Invariant}\par\vspace{4pt}

{\large\itshape\raggedright
  Cross-Precision Output Divergence in Large Language Model Inference}\par\vspace{6pt}

{\small\textbf{By} Gaoyuan Du, Anam Nawaz Khan, Rex Zhou, Xiaoyang Liu,
  Deepayan Chakrabarti, Fnu Suya \& Xueping Li\\
  (UTK \textbullet{} University of Chicago \textbullet{} Amazon \textbullet{} UT Austin)}\par\vspace{2pt}

{\small\color{keywordgray}arXiv:2609.26621 \textbullet{} September 2026}
\par\vspace{2pt}\noindent{\color{rulered}\rule{\columnwidth}{0.6pt}}\vspace{6pt}

BF16 and FP16 greedy decoding from the same model on the same hardware
diverges for 49--100\% of prompts across six models and three benchmarks;
the divergence is predicted not by deep-layer error accumulation but by
the top-two logit margin at the language model head, and selective
FP32 recomputation of that head at tight-margin steps recovers
+22--36\,pp exact agreement at only $\pm$4\% latency overhead.

\begin{mdframed}[backgroundcolor=boxgray,linecolor=rulered,linewidth=1pt,
    innerleftmargin=6pt,innerrightmargin=6pt,
    innertopmargin=6pt,innerbottommargin=6pt]
\textbf{\small AT A GLANCE}\par\vspace{4pt}
\begin{description}[leftmargin=0pt,labelsep=4pt,itemsep=2pt,parsep=0pt,
    font=\small\bfseries]
  \item[Problem:] Greedy decoding is assumed to be deterministic given a
    fixed model and prompt, but BF16 and FP16 produce different outputs
    on identical hardware at a surprisingly high rate.
  \item[Why it matters:] Cross-precision divergence confounds evaluations,
    complicates deployment, and undermines reproducibility guarantees.
  \item[Method:] Empirical error-propagation analysis traces divergence to
    the top-two logit margin at the LM head; selective FP32 recomputation
    at low-margin steps corrects most divergences.
  \item[Result:] +22--36\,pp exact agreement recovered on A10G at $\pm$4\%
    latency overhead; approach validated across three GPU architectures.
  \item[Evidence:] Six models (1.1B--7B, four families), three benchmarks
    (MMLU, TriviaQA, GSM8K).
\end{description}
\end{mdframed}

\section{The Big Picture}

Large language model evaluations routinely compare numbers across
frameworks, hardware generations, and model serving stacks. An implicit
assumption underlies all such comparisons: greedy decoding is
\emph{deterministic} given a fixed model. Feed in the same model weights,
the same prompt, and the same hardware, and you get the same output.

This paper breaks that assumption. It shows that switching between BF16
and FP16---two floating-point formats in equally widespread use---produces
different greedy outputs from the same model on the same GPU for 49 to
100\% of prompts, depending on the model and benchmark. The divergence is
not marginal: a single token flip typically cascades into a fully different
response trajectory.

The culprit is not the 22 layers of accumulated numerical error that one
might expect. It is a \emph{local} phenomenon: whether the top-two logits
at the language model head are separated by more than the cross-precision
perturbation at that point in the computation.

\section{Why Existing Approaches Fall Short}

\textbf{Same-precision nondeterminism} (parallelism-order variation within
BF16 or FP16) is well-studied and mitigated by deterministic CUDA kernels.
It does not explain \emph{cross-precision} divergence, which is systematic
and reproducible.

\textbf{Full FP32 inference} eliminates the problem but doubles memory usage
and costs roughly 1.8$\times$ in latency---impractical for production.

\textbf{Precision-agnostic evaluation} (reporting numbers without specifying
precision) silently inflates or deflates benchmark scores by comparison-format
mismatches.

No targeted, low-cost mitigation existed before this work.

\section{Core Mechanism}

\textbf{Systematic Characterization.} The authors evaluate six LLMs
(1.1B--7B, four families) with greedy decoding in BF16 and FP16 on MMLU,
TriviaQA, and GSM8K on identical A10G hardware. For each prompt, they record
whether the full trajectory diverges.

\textbf{Error-Propagation Analysis.} For each diverging prompt, they trace
the first-flip location and measure accumulated numerical error at each layer
and at the LM head. The finding: 22 layers of body error do not differentiate
flipping from non-flipping steps. Divergence is determined at the LM head.

\textbf{Selective FP32 Head Recomputation.} When the top-two logit margin
at the LM head falls below a threshold $\tau$, the head matrix multiplication
is re-run at FP32 for that decoding step. Steps with large margins are left
at native precision.

\section{Technical Formulation}

Let $\mathbf{h} \in \mathbb{R}^d$ be the last-layer hidden state and
$\mathbf{W} \in \mathbb{R}^{V \times d}$ the LM head weight matrix.
Logits under precision $p \in \{\mathrm{bf16}, \mathrm{fp16}\}$:
\[
\mathbf{z}^{(p)} = \mathbf{W}^{(p)} \mathbf{h}^{(p)} + \boldsymbol{\epsilon}^{(p)}
\]
where $\boldsymbol{\epsilon}^{(p)}$ is accumulated rounding error.

A token flip occurs when:
\[
\arg\max \mathbf{z}^{(\mathrm{bf16})} \neq \arg\max \mathbf{z}^{(\mathrm{fp16})}
\]

Empirically, the flip probability is well-predicted by:
\[
\Pr[\mathrm{flip}] \approx
  \Pr\!\left[\Delta_{\mathrm{top2}} <
    \bigl\|\boldsymbol{\epsilon}^{(\mathrm{bf16})} - \boldsymbol{\epsilon}^{(\mathrm{fp16})}
    \bigr\|_{\mathrm{top2\text{-}dir}}\right]
\]
where $\Delta_{\mathrm{top2}} = z_{(1)} - z_{(2)}$ is the top-two margin.

\textbf{Selective FP32 trigger:}
\[
\text{recompute at FP32} \iff \Delta_{\mathrm{top2}} < \tau
\]

\section{Learning or Inference Procedure}

No training is required. The mitigation applies at each decoding step:
\begin{enumerate}[itemsep=2pt,parsep=0pt]
  \item Run the transformer forward pass in native precision (BF16 or FP16).
  \item Compute logits and the top-two margin $\Delta_{\mathrm{top2}}$.
  \item If $\Delta_{\mathrm{top2}} < \tau$: re-run the LM head matmul in FP32;
    use FP32 logits for this step.
  \item Select the greedy token; proceed.
\end{enumerate}
The threshold $\tau$ is calibrated per hardware on a small validation set.

\section{What the Guarantee Says}

Selective FP32 head recomputation is an empirical, not worst-case, mitigation.
It cannot guarantee identical outputs to any reference because cross-precision
error accumulates throughout the forward pass, not only at the head. However,
because the head is the single decision point and body error is shown not to
be predictive, covering tight-margin head steps recovers most of the divergence.
The remaining divergence (about 10--30\%) stems from body-precision differences
that the head-only fix does not address.

\section{Experimental Findings}

\begin{itemize}[itemsep=2pt,parsep=0pt]
  \item \textbf{Divergence rate:} 49--100\% of prompts across 6 models
    and 3 benchmarks; rate varies by model family and task type.
  \item \textbf{Cascade effect:} A single first-token flip cascades into
    full trajectory divergence in $>$80\% of cases.
  \item \textbf{Error-propagation:} Accumulated body error over 22 layers
    does not predict flipping; top-two margin at the LM head is the dominant
    predictor.
  \item \textbf{Selective FP32 on A10G:} +22--36\,pp exact-agreement
    recovery at $\pm$4\% latency overhead.
  \item \textbf{L4 and A100:} +12--21\,pp recovery (smaller because
    baseline divergence rates differ by GPU architecture).
\end{itemize}

\section{Ablations and Interpretation}

\textbf{Threshold sensitivity:} Agreement recovery is robust across a 5$\times$
range of $\tau$; the optimal value balances coverage (catching flip-prone steps)
against unnecessary FP32 overhead.

\textbf{Full FP32:} Achieves near-zero divergence but at 2$\times$ memory and
$\sim$1.8$\times$ latency---impractical for production serving.

\textbf{Body recomputation:} Re-running the full transformer in FP32 at
flip-predicted steps adds disproportionate overhead for marginal additional gain.

\textbf{Vocabulary size effect:} Larger vocabularies have smaller average
top-two margins (more competing candidates per step) and higher divergence rates,
suggesting the problem worsens as vocabulary size grows.

\textbf{Task type:} Open-ended generation (TriviaQA) shows higher divergence
than constrained generation (MMLU multiple-choice), as cascading divergence
has more opportunity to propagate in longer unconstrained outputs.

\section*{Reference}

Gaoyuan Du, Anam Nawaz Khan, Rex Zhou, Xiaoyang Liu, Deepayan Chakrabarti,
Fnu Suya, Xueping Li.
\textbf{Greedy Decoding Is Not Precision-Invariant: Cross-Precision Output
Divergence in LLM Inference}.
arXiv:2609.26621, September 2026.\\
\url{https://arxiv.org/abs/2609.26621}

\end{document}
